\scanpage{44}
Ici, prenant les top. étales sur les préschémas, et les faisceaux
\unsure{cts} $A_X$ $(\mathbb Z/n\mathbb Z=A)$,
tous les morphismes de topos annelés envisagés seront plats.

\figure{\begin{tikzpicture}
\draw (0,0)--(1,0);\draw (.18,-.12)--(.77,-.12);
\end{tikzpicture}}

\margin{$\unsure{F}$}
\diagram{\begin{tikzcd}[column sep=large,row sep=large]
X\arrow[d,dash,"f"']&X'\arrow[l,"h"']\arrow[d,"f'"]\\
Y&Y'\arrow[l,"g"]
\end{tikzcd}}
Si $f$ quelconque, on définit pour $F^\bullet\in\mathrm{Ob}D^+(X)$
un hom (dans la catégorie dérivée $D^+(Y')$)
\[
(*)\quad \deleted{Rf'_*(h^*F^\bullet)}
\deleted{\longleftarrow}
\boxed{g^*(Rf_*(F^\bullet))\longrightarrow Rf'_*(\deleted{\unsure{g}}h^*(F^\bullet))}
\]
Passant aux faisceaux de coh., cela redonne
\[
(*_i)\quad g^+R^i(f_*(F))\longrightarrow R^if'_*(h^*(F)).
\]
\uline{Lorsque} $f$ \uline{est propre} ce est un Isom,
puisqu'on sait que les $(*_i)$ le sont.
\mspage{43}
